\documentclass[12pt]{article}
\usepackage[utf8]{inputenc}
\usepackage[T1]{fontenc}
\usepackage{hyperref}
\title{Terik Language Resources for K-12 Curriculum Development}
\author{Compiled by Claude Sonnet 5 (claude-sonnet-5) for Andr\'as Kornai}
\date{}
\begin{document}
\maketitle
\section{Introduction}
Terik (Nyang'ori; ISO 639-3 tec) is a Southern Nilotic language of the Kalenjin group in western Kenya, mutually intelligible with Nandi, with about 300,000 reported speakers, mostly middle-aged or older. Search found little beyond overview pages.
\section{Primary Educational Resources}
None found. No Terik primers or textbooks appeared in search.
\section{Linguistic and Reference Materials}
Overview: \url{https://en.wikipedia.org/wiki/Terik_language} and \url{https://en.wikipedia.org/wiki/Terik_people}. Omniglot has an alphabet page: \url{https://www.omniglot.com/writing/terik.htm}. The Wikipedia article cites Bernd Heine's 1992 chapter ``Dialect death: the case of Terik'' (not found online in search; no URL).
\section{Religious and Community Resources}
None found.
\section{Digital and Audio Resources}
None found.
\section{Curriculum Development Context}
Documented Terik material is minimal; a curriculum would require new documentation, possibly leveraging the better-resourced related Kalenjin/Nandi varieties (not searched here).
\end{document}